%% =============================================================================
%% CPS slot for the final. Replaces Binary Decision Trees (too hard; its
%% `compress` refsol also hit the value restriction -- both copies kept at
%% q/bdt.tex and q/bdt_inlined.tex) and then "Short Circuits" (which asked
%% students to invent discarding the continuation -- lecture13's `match` does it
%% at `Zero => false`, but the lectures never name or explain the idea).
%%
%% This question uses ONLY the idiom lecture13 develops explicitly: the
%% continuation receives the UNCONSUMED REMAINDER of the input, and nested calls
%% thread it through --
%%     Times (r1, r2) => match r1 cs (fn cs' => match r2 cs' k)
%% Framed as parsing rather than tree-vs-sequence comparison, so it does not
%% overlap the practice exam's samefringe (streams) or its fillAll (the
%% list-concatenating `@` continuation).
%%
%% Verified in SML/NJ: all three sample trees round-trip through
%% deserialize o serialize; trailing junk and truncated input both raise
%% Malformed.
%% =============================================================================

\titledquestion{Serial Offenders}

\textbf{Serial Offenders}

We want to write a tree to a flat sequence of tokens and read it back. Our trees
carry no data --- only shape:
\begin{codeblock}
  datatype tree  = Leaf | Node of tree * tree
  datatype token = LeafTok | NodeTok

  exception Malformed
\end{codeblock}

A tree is written in \textbf{preorder}: a \code{Leaf} becomes the single token
\code{LeafTok}, and a \code{Node} becomes \code{NodeTok} followed by the tokens
of its left subtree and then those of its right. For example, the tree
\code{Node (Node (Leaf, Leaf), Leaf)} is written as
\begin{codeblock}
  [NodeTok, NodeTok, LeafTok, LeafTok, LeafTok]
\end{codeblock}

Notice that there are no parentheses or separators in the output. Reading it back
works only because each \code{NodeTok} promises that exactly two more trees
follow.

\begin{parts}

  \part[5]\hfill

    As a warm-up, implement \code{serialize}.

    \spec
      {serialize}
      {\code{tree -> token list}}
      {\code{true}}
      {\code{serialize T} $\eeq$ the preorder token list of \code{T}, as
      described above}

    \begin{solutionorbox}[12em]
      \begin{codeblock}
        fun serialize Leaf = [LeafTok]
          | serialize (Node (l, r)) =
              NodeTok :: serialize l @ serialize r
      \end{codeblock}
    \end{solutionorbox}

  \newpage

  \part[12]\hfill

    Now we read one back. The difficulty is that when we parse the left subtree
    of a \code{Node}, we do not know in advance how many tokens it will consume
    --- so the left subtree must tell us where it stopped.

    We handle this in CPS: the continuation is given both the tree that was
    parsed \textbf{and the tokens that are left over}.

    \spec
      {parse}
      {\\ {\small\code{token list -> (tree * token list -> 'a) -> 'a}}}
      {\code{k} is total}
      {\code{parse ts k} $\eeq$ \code{k (T, ts')}, where \code{T} is the tree
      described by a prefix of \code{ts} and \code{ts'} is the remaining tokens.
      If \code{ts} does not begin with a well-formed tree, \code{parse ts k}
      raises \code{Malformed}.}

    \begin{solutionorbox}[20em]
      \begin{codeblock}
        fun parse (LeafTok :: ts) k = k (Leaf, ts)
          | parse (NodeTok :: ts) k =
              parse ts (fn (lt, ts') =>
                parse ts' (fn (rt, ts'') =>
                  k (Node (lt, rt), ts'')))
          | parse [] k = raise Malformed
      \end{codeblock}

      A \code{LeafTok} consumes exactly one token, so the continuation gets
      \code{Leaf} and everything after it.

      For a \code{NodeTok} we parse the left subtree from the tokens that follow.
      Its continuation receives \code{lt} together with \code{ts'}, the tokens
      the left subtree did not use --- and \textit{those} are what the right
      subtree parses from. The innermost continuation then has both subtrees and
      the final leftovers \code{ts''}, so it can build the node and hand it on.

      Threading the remainder through the continuations is what removes the need
      to know any subtree's length in advance.
    \end{solutionorbox}

  \newpage

  \part[5]\hfill

    Finally, use \code{parse} to implement \code{deserialize}, which must consume
    the \textbf{entire} token list.

    \spec
      {deserialize}
      {\code{token list -> tree}}
      {\code{true}}
      {\code{deserialize ts} $\eeq$ \code{T} if \code{ts} $\eeq$
      \code{serialize T}, and raises \code{Malformed} otherwise}

    \begin{solutionorbox}[14em]
      \begin{codeblock}
        fun deserialize ts =
          parse ts (fn (t, rest) =>
            if null rest then t
            else raise Malformed)
      \end{codeblock}

      \code{parse} is happy to stop early --- it reports what it consumed and
      what it did not. So the continuation is where ``and nothing may be left
      over'' is enforced: it accepts the tree only when \code{rest} is empty.

      This is the useful part of returning the remainder rather than hiding it.
      The same \code{parse} would serve a caller that wanted to read a tree and
      then keep going with the rest of the tokens; only the continuation would
      differ.
    \end{solutionorbox}

\end{parts}
